
\uline{où} @@M0@@ \uline{est la
classe de Chern totale} (V 6.7).

Ce résultat se montre en recopiant la démonstration de 4.6, et en utilisant 5.6 (comparer [2], Exp 4).

\textbf{Corollaire 5.8.} \uline{Avec les notations et les hypothèses de} 5.1, \uline{on a pour} \uline{tout} @@M1@@ :
@@M2@@

Conséquence immédiate de 5.3 et de la formule de projection.

\textbf{Proposition 5.9.} \uline{Avec les notations et les hypothèses de} 5.1, \uline{on a pour} \uline{tout} @@M3@@ :
@@M4@@
\uline{En particulier}, \uline{si on pose} @@M5@@, \uline{on obtient} :
@@M6@@

On définit une application
@@M7@@
en posant :
@@M8@@
Ce n'est pas en général un homomorphisme de groupes; néanmoins, la relation suivante est vérifiée pour @@M9@@, @@M10@@ :
@@M11@@
En effet, elle se vérifie coefficient par coefficient, et résulte immédiatement de la formule de projection.

De la relation @@M12@@, on déduit donc :
